% ECONOMICS_PROBLEM: {"id": "L26-436", "title": "Declining business dynamism", "jel_code": "L26", "source_id": "OP-0265"}
% !TeX program = xelatex
\documentclass[11pt,a4paper]{article}
\usepackage[margin=23mm]{geometry}
\usepackage{fontspec}
\setmainfont{Latin Modern Roman}
\usepackage{amsmath,amssymb,mathtools}
\usepackage{xcolor,enumitem,booktabs,longtable,array,xurl,hyperref}
\definecolor{muted}{HTML}{53636C}
\hypersetup{colorlinks=true,urlcolor=blue,linkcolor=blue}
\setlength{\parindent}{0pt}
\setlength{\parskip}{5pt}
\newcommand{\R}{\mathbb R}
\newcommand{\E}{\mathbb E}
\newcommand{\N}{\mathbb N}
\newcommand{\ind}{\mathbf 1}
\newcommand{\Var}{\operatorname{Var}}
\newcommand{\Cov}{\operatorname{Cov}}
\newcommand{\supp}{\operatorname{supp}}
\DeclareMathOperator*{\argmin}{arg\,min}
\DeclareMathOperator*{\argmax}{arg\,max}
\newcommand{\entrymeta}[3]{{\sffamily\small\color{muted}\textbf{\texttt{#1}}\quad|\quad #2\par #3\par}}
\newcommand{\Problemref}[1]{\texttt{#1}}
\newcommand{\JELref}[2]{\texttt{#2} (JEL \texttt{#1})}
\begin{document}
\section*{Declining business dynamism}
\textbf{Problem ID:} \texttt{L26-436}\quad \textbf{JEL:} \texttt{L26}\par
% BEGIN ECONOMICS STATEMENT
\label{problem:OP-0265}
{\sffamily\small\color{muted}Original subject group: Growth and productivity\par}
\label{definitions:problem:30007546}
\textbf{Audit classification:} Explicit published open question. Reconciling aggregate and industry patterns is an explicit research gap; the four-block decomposition is editorial. Verification concerns the cited version, not first priority or proof that this exact editorial specification remains unsolved on October 8, 2026.
\subsection*{Definitions and precise research target}
\paragraph{Editorial mathematical specification.} Let \(D_t\) be a predeclared business-dynamism index built from firm entry, exit, and employment reallocation in a fixed population. Consider a structural firm-dynamics model with parameter blocks \(\theta=(\theta^N,\theta^C,\theta^I,\theta^R)\), representing demographics and labor supply, competitive conditions, intangible technology and diffusion, and regulation. Specify how each block affects entry costs, incumbent innovation, and firm growth; a residual trend must not silently receive a causal label.
For initial and final parameter vectors \(\theta_0,\theta_1\), let \(D(\theta)\) be the equilibrium index. Decompose \(D(\theta_1)-D(\theta_0)\) into block contributions using a declared Shapley ordering average over counterfactual combinations. Identify those contributions from moments containing both aggregate trends and cross-industry variation, augmented by credible policy or technological instruments. In particular, test whether the same parameter changes explain industry-level relationships among markups, entry, and reallocation rather than only matching aggregate comovement.
Report observationally equivalent decompositions and parameter uncertainty, distinguish measured markups from primitive market power, and assess robustness to firm-population coverage. A successful result reconciles time-series decline with cross-sectional evidence and predicts held-out changes. The target is a quantitatively disciplined explanation for a stated period and economy, not a presumption that a single driver must dominate universally.

One precise index is \(D_t=\omega_1 E_t+\omega_2 X_t+\omega_3J_t\), where \(E_t\) and \(X_t\) are firm entry/exit fractions in a fixed covered population, \(J_t\) is a predeclared employment-reallocation rate, and nonnegative fixed weights sum to one. Do not mix establishment entry with firm entry without a mapping. For the four parameter blocks \(K\), define \(v(S)=D(\theta_1^S,\theta_0^{K\setminus S})\); the Shapley contribution for each block uses the same factorial weights as standard finite-set allocation. Every mixed-parameter economy must have a defined selected equilibrium; otherwise that decomposition is not defined. The allocation is convention-dependent, while identification requires data restrictions distinguishing parameter blocks. A measured markup is a price-to-estimated-marginal-cost ratio and does not by itself identify the underlying competitive primitive.
\subsection*{Variants and assumption changes}
The following are \emph{editorial variants}, not additional published open conjectures.
\begin{enumerate}
\item \emph{Population mechanism.} Vary labor-force growth and age composition while holding entry technology and competition parameters fixed. Require the resulting model to match age-specific firm entry and industry reallocation, beyond the aggregate decline.
\item \emph{Industry reconciliation.} Estimate a competition/entry mechanism that predicts both aggregate time trends and cross-industry markup-entry patterns. Compare alternative markup measures and public/private firm coverage under the same model restrictions.
\end{enumerate}
\subsection*{References and source audit}
\begin{itemize}
\item Brian C. Albrecht; Ryan A. Decker. \emph{Rising Markups and Declining Business Dynamism: Evidence From the Industry Cross Section}. 2024. Locator: Section2, Related literature, final paragraph on reconciliation; dated March 8, 2024. Primary text checked. \url{https://www.federalreserve.gov/econres/notes/feds-notes/rising-markups-and-declining-business-dynamism-evidence-from-the-industry-cross-section-20240308.html}.
\end{itemize}
Reconciling aggregate and industry patterns is an explicit research gap; the four-block decomposition is editorial.
\begin{enumerate}\raggedright
\item\hypertarget{definitions:source:30007546:1}{} Brian C. Albrecht; Ryan A. Decker. 2024. \emph{Rising Markups and Declining Business Dynamism: Evidence From the Industry Cross Section}. Section2, Related literature, final paragraph on reconciliation; dated March8,2024.
\url{https://www.federalreserve.gov/econres/notes/feds-notes/rising-markups-and-declining-business-dynamism-evidence-from-the-industry-cross-section-20240308.html}.
\end{enumerate}

\subsection*{Common conventions: Source mathematical conventions}

These conventions apply to each entry unless explicitly replaced.

\textbf{Models and identification.} A primitive model \(m\) specifies all probability laws, feasible choices, information, equilibrium and measurement rules used by its target. The admissible class \(\mathcal M\) is fixed independently of outcomes. If \(P_O(m)\) is its observable probability law and \(\tau(m)\) the target, its identified set at observable law \(P\) is
\[
 \mathcal I_\tau(P)=\{\tau(m):m\in\mathcal M,\ P_O(m)=P\}.
\]
Point identification means that this set is a singleton. An empty set signals incompatibility of data and assumptions. Coordinate bounds need not describe the joint set. Bounds are sharp only if every retained target is attainable or arbitrarily approachable by compatible models. Finite bounds require explicit restrictions on unobserved quantities; unrestricted confounding, hidden wealth, extrapolation or arbitrary equilibrium selection can yield uninformative sets.

\textbf{Causal interventions.} A potential outcome \(Y_i(p)\) is the outcome under a complete policy regime \(p\), including timing, financing, information and market spillovers. The target population and units are fixed before treatment. With interference, use the complete assignment vector or a justified exposure mapping; individual no-interference notation alone cannot identify an economy-wide intervention. A difference of mean potential outcomes does not require the cross-world joint distribution. Conditioning on post-treatment migration, survival, enrollment or employment changes the estimand and needs separate justification.

\textbf{Statistical statements.} An estimator, confidence set, forecast or test specifies a sampling law, model class and loss/coverage criterion. Uniform \((1-\alpha)\)-coverage means \(\inf_{m\in\mathcal M}\Pr_m\{\tau(m)\in C_n\}\geq1-\alpha\), or an explicitly asymptotic version. The whole parameter space trivially has coverage, so informativeness requires a stated width, risk or power objective. A forecasting improvement is relative to a named benchmark, information set and evaluation population.

\textbf{Equilibrium and optimization.} An equilibrium comprises feasible plans, optimality, beliefs where needed, budget constraints and all market/resource clearing. The primitives or a stated benchmark define each correspondence \(\mathcal E(p;m)\). Nonemptiness is assumed or proved, never inferred just from compact policy sets. A selected-equilibrium target uses a declared selection rule; a robust target quantifies over all permitted equilibria. Use suprema or infima unless attainment is established. An \(\varepsilon\)-optimal policy is feasible and within \(\varepsilon\) of the finite optimum value. Report an empty feasible set separately.

\textbf{Welfare and accounting.} Normative utility, interpersonal normalization, social weights, numeraire, discounting and the use of public receipts are primitives. Behavioral utility need not coincide with the welfare criterion. Household welfare includes ownership income and rents once; adding profits again to compensating variation generally double-counts them. A monetary equivalent is defined by an explicit transfer and price convention, with monotonicity and range conditions. Average effects, mechanism contributions, transfers and welfare losses are different targets. Infinite sums require convergence and dynamic feasible plans require terminal or no-Ponzi conditions.

\textbf{Source versions.} A working-paper date, revision date and journal publication year are distinguished. A DOI for a journal version is not automatically the DOI of an earlier manuscript. Locators refer to the stated version and printed pages where specified. A metadata-only check does not verify a claim about a full-text conclusion. Every entry's source subsection narrows the provenance claim accordingly.
% END ECONOMICS STATEMENT
\end{document}
